Il Cloud Computing è un paradigma oramai ampiamente affermato in quanto offre la possibilità di accedere a diversi tipi di risorse hardware e software mediante un paradigma basato su servizi. L'espansione del Cloud ha spinto numerose organizzazioni a diventare Cloud Provider, ovvero fornitori di servizi. Tale fenomeno ha spinto la nascita del multi-cloud. Il multi-cloud prevede l'utilizzo all'interno di una singola architettura di molteplici Cloud provider garantendo una maggiore affidabilità, un costo più vantaggioso e servizi più eterogenei.
Tuttavia l'utilizzo del multi-cloud non è banale in quanto ogni Cloud provider mette a disposizione strumenti e tecnologie anche molto diverse tra di loro che rendono complesso il loro utilizzo e incrementano il tempo per l'apprendimento.
Tale problematica trova un particolare interesse all'interno della comunità scientifica che spinge la ricerca e lo sviluppo di strumenti che permettano di rendere più semplice ed immediato l'utilizzo del Cloud Computing. In questo contesto nasce FLY, un Domain-Specific-Language per il Calcolo Scientifico sul multi-cloud che ha l’obiettivo di semplificare lo sviluppo di applicazioni che sfruttino la potenza computazionale offerta da molteplici Cloud provider. FLY sfrutta il paradigma Function-as-a-Service mediante il costrutto nativo FLY function che permette di ottenere alta scalabilità e alte prestazioni. Il maggiore vantaggio di FLY è permettere a qualsiasi utente di sfruttare a pieno le potenzialità del multi-cloud astraendo l'interazione con il provider che viene totalmente gestita dal linguaggio.
Oggigiorno, lavorare con i dati è una necessità in qualsiasi contesto e applicazione reale e negli ultimi anni, la mole di dati da gestire e la velocità con cui vengono prodotti è aumentata sempre di più. In aggiunta a ciò, l'eterogeneità delle fonti da cui provengono tali dati introduce una serie di requisiti aggiuntivi nei sistemi di gestione delle informazioni come alte prestazioni, bassa latenza e elevata flessibilità. 
La nascita dei database NoSQL viene incontro a questo tipo di necessità fornendo schemi flessibili per la memorizzazione e la gestione dei dati e la capacità di scalare velocemente sfruttando cluster e sistemi distribuiti.  Le caratteristiche dei database NoSQL li rendono quindi particolarmente adatti ad applicazioni che richiedono alta potenza di calcolo e ad essere implementati sfruttando il Cloud Computing, e per questo motivo è risultato necessario implementare un supporto per i database NoSQL all’interno di FLY.
Questo lavoro di tesi ha lo scopo di presentare l'introduzione all'interno del linguaggio FLY del supporto ai database NoSQL.